\IfFileExists{stacks-project.cls}{%
\documentclass{stacks-project}
}{%
\documentclass{amsart}
}

% For dealing with references we use the comment environment
\usepackage{verbatim}
\newenvironment{reference}{\comment}{\endcomment}
%\newenvironment{reference}{}{}
\newenvironment{slogan}{\comment}{\endcomment}
\newenvironment{history}{\comment}{\endcomment}

% For commutative diagrams we use Xy-pic
\usepackage[all]{xy}

% We use 2cell for 2-commutative diagrams.
\xyoption{2cell}
\UseAllTwocells

% We use multicol for the list of chapters between chapters
\usepackage{multicol}

% This is generally recommended for better output
\usepackage{lmodern}
\usepackage[T1]{fontenc}

% For cross-file-references
\usepackage{xr-hyper}

% Package for hypertext links:
\usepackage{hyperref}

% For any local file, say "hello.tex" you want to link to please
% use \externaldocument[hello-]{hello}
\externaldocument[introduction-]{introduction}
\externaldocument[conventions-]{conventions}
\externaldocument[sets-]{sets}
\externaldocument[categories-]{categories}
\externaldocument[topology-]{topology}
\externaldocument[sheaves-]{sheaves}
\externaldocument[sites-]{sites}
\externaldocument[stacks-]{stacks}
\externaldocument[fields-]{fields}
\externaldocument[algebra-]{algebra}
\externaldocument[brauer-]{brauer}
\externaldocument[homology-]{homology}
\externaldocument[derived-]{derived}
\externaldocument[simplicial-]{simplicial}
\externaldocument[more-algebra-]{more-algebra}
\externaldocument[smoothing-]{smoothing}
\externaldocument[modules-]{modules}
\externaldocument[sites-modules-]{sites-modules}
\externaldocument[injectives-]{injectives}
\externaldocument[cohomology-]{cohomology}
\externaldocument[sites-cohomology-]{sites-cohomology}
\externaldocument[dga-]{dga}
\externaldocument[dpa-]{dpa}
\externaldocument[sdga-]{sdga}
\externaldocument[hypercovering-]{hypercovering}
\externaldocument[schemes-]{schemes}
\externaldocument[constructions-]{constructions}
\externaldocument[properties-]{properties}
\externaldocument[morphisms-]{morphisms}
\externaldocument[coherent-]{coherent}
\externaldocument[divisors-]{divisors}
\externaldocument[limits-]{limits}
\externaldocument[varieties-]{varieties}
\externaldocument[topologies-]{topologies}
\externaldocument[descent-]{descent}
\externaldocument[perfect-]{perfect}
\externaldocument[more-morphisms-]{more-morphisms}
\externaldocument[flat-]{flat}
\externaldocument[groupoids-]{groupoids}
\externaldocument[more-groupoids-]{more-groupoids}
\externaldocument[etale-]{etale}
\externaldocument[chow-]{chow}
\externaldocument[intersection-]{intersection}
\externaldocument[pic-]{pic}
\externaldocument[weil-]{weil}
\externaldocument[adequate-]{adequate}
\externaldocument[dualizing-]{dualizing}
\externaldocument[duality-]{duality}
\externaldocument[discriminant-]{discriminant}
\externaldocument[derham-]{derham}
\externaldocument[local-cohomology-]{local-cohomology}
\externaldocument[algebraization-]{algebraization}
\externaldocument[curves-]{curves}
\externaldocument[resolve-]{resolve}
\externaldocument[models-]{models}
\externaldocument[functors-]{functors}
\externaldocument[equiv-]{equiv}
\externaldocument[pione-]{pione}
\externaldocument[etale-cohomology-]{etale-cohomology}
\externaldocument[proetale-]{proetale}
\externaldocument[relative-cycles-]{relative-cycles}
\externaldocument[more-etale-]{more-etale}
\externaldocument[trace-]{trace}
\externaldocument[crystalline-]{crystalline}
\externaldocument[spaces-]{spaces}
\externaldocument[spaces-properties-]{spaces-properties}
\externaldocument[spaces-morphisms-]{spaces-morphisms}
\externaldocument[decent-spaces-]{decent-spaces}
\externaldocument[spaces-cohomology-]{spaces-cohomology}
\externaldocument[spaces-limits-]{spaces-limits}
\externaldocument[spaces-divisors-]{spaces-divisors}
\externaldocument[spaces-over-fields-]{spaces-over-fields}
\externaldocument[spaces-topologies-]{spaces-topologies}
\externaldocument[spaces-descent-]{spaces-descent}
\externaldocument[spaces-perfect-]{spaces-perfect}
\externaldocument[spaces-more-morphisms-]{spaces-more-morphisms}
\externaldocument[spaces-flat-]{spaces-flat}
\externaldocument[spaces-groupoids-]{spaces-groupoids}
\externaldocument[spaces-more-groupoids-]{spaces-more-groupoids}
\externaldocument[bootstrap-]{bootstrap}
\externaldocument[spaces-pushouts-]{spaces-pushouts}
\externaldocument[spaces-chow-]{spaces-chow}
\externaldocument[groupoids-quotients-]{groupoids-quotients}
\externaldocument[spaces-more-cohomology-]{spaces-more-cohomology}
\externaldocument[spaces-simplicial-]{spaces-simplicial}
\externaldocument[spaces-duality-]{spaces-duality}
\externaldocument[formal-spaces-]{formal-spaces}
\externaldocument[restricted-]{restricted}
\externaldocument[spaces-resolve-]{spaces-resolve}
\externaldocument[formal-defos-]{formal-defos}
\externaldocument[defos-]{defos}
\externaldocument[cotangent-]{cotangent}
\externaldocument[examples-defos-]{examples-defos}
\externaldocument[algebraic-]{algebraic}
\externaldocument[examples-stacks-]{examples-stacks}
\externaldocument[stacks-sheaves-]{stacks-sheaves}
\externaldocument[criteria-]{criteria}
\externaldocument[artin-]{artin}
\externaldocument[quot-]{quot}
\externaldocument[stacks-properties-]{stacks-properties}
\externaldocument[stacks-morphisms-]{stacks-morphisms}
\externaldocument[stacks-limits-]{stacks-limits}
\externaldocument[stacks-cohomology-]{stacks-cohomology}
\externaldocument[stacks-perfect-]{stacks-perfect}
\externaldocument[stacks-introduction-]{stacks-introduction}
\externaldocument[stacks-more-morphisms-]{stacks-more-morphisms}
\externaldocument[stacks-geometry-]{stacks-geometry}
\externaldocument[moduli-]{moduli}
\externaldocument[moduli-curves-]{moduli-curves}
\externaldocument[examples-]{examples}
\externaldocument[exercises-]{exercises}
\externaldocument[guide-]{guide}
\externaldocument[desirables-]{desirables}
\externaldocument[coding-]{coding}
\externaldocument[obsolete-]{obsolete}
\externaldocument[fdl-]{fdl}
\externaldocument[index-]{index}

% Theorem environments.
%
\theoremstyle{plain}
\newtheorem{theorem}[subsection]{Theorem}
\newtheorem{proposition}[subsection]{Proposition}
\newtheorem{lemma}[subsection]{Lemma}

\theoremstyle{definition}
\newtheorem{definition}[subsection]{Definition}
\newtheorem{example}[subsection]{Example}
\newtheorem{exercise}[subsection]{Exercise}
\newtheorem{situation}[subsection]{Situation}

\theoremstyle{remark}
\newtheorem{remark}[subsection]{Remark}
\newtheorem{remarks}[subsection]{Remarks}

\numberwithin{equation}{subsection}

% Macros
%
\def\lim{\mathop{\mathrm{lim}}\nolimits}
\def\colim{\mathop{\mathrm{colim}}\nolimits}
\def\Spec{\mathop{\mathrm{Spec}}}
\def\Hom{\mathop{\mathrm{Hom}}\nolimits}
\def\Ext{\mathop{\mathrm{Ext}}\nolimits}
\def\SheafHom{\mathop{\mathcal{H}\!\mathit{om}}\nolimits}
\def\SheafExt{\mathop{\mathcal{E}\!\mathit{xt}}\nolimits}
\def\Sch{\mathit{Sch}}
\def\Mor{\mathop{\mathrm{Mor}}\nolimits}
\def\Ob{\mathop{\mathrm{Ob}}\nolimits}
\def\Sh{\mathop{\mathit{Sh}}\nolimits}
\def\NL{\mathop{N\!L}\nolimits}
\def\CH{\mathop{\mathrm{CH}}\nolimits}
\def\proetale{{pro\text{-}\acute{e}tale}}
\def\etale{{\acute{e}tale}}
\def\QCoh{\mathit{QCoh}}
\def\Ker{\mathop{\mathrm{Ker}}}
\def\Im{\mathop{\mathrm{Im}}}
\def\Coker{\mathop{\mathrm{Coker}}}
\def\Coim{\mathop{\mathrm{Coim}}}

% Boxtimes
%
\DeclareMathSymbol{\boxtimes}{\mathbin}{AMSa}{"02}

%
% Macros for moduli stacks/spaces
%
\def\QCohstack{\mathcal{QC}\!\mathit{oh}}
\def\Cohstack{\mathcal{C}\!\mathit{oh}}
\def\Spacesstack{\mathcal{S}\!\mathit{paces}}
\def\Quotfunctor{\mathrm{Quot}}
\def\Hilbfunctor{\mathrm{Hilb}}
\def\Curvesstack{\mathcal{C}\!\mathit{urves}}
\def\Polarizedstack{\mathcal{P}\!\mathit{olarized}}
\def\Complexesstack{\mathcal{C}\!\mathit{omplexes}}
% \Pic is the operator that assigns to X its picard group, usage \Pic(X)
% \Picardstack_{X/B} denotes the Picard stack of X over B
% \Picardfunctor_{X/B} denotes the Picard functor of X over B
\def\Pic{\mathop{\mathrm{Pic}}\nolimits}
\def\Picardstack{\mathcal{P}\!\mathit{ic}}
\def\Picardfunctor{\mathrm{Pic}}
\def\Deformationcategory{\mathcal{D}\!\mathit{ef}}
\usepackage{cleveref}
\crefname{originalreference}{original reference}{original references}
\Crefname{originalreference}{Original reference}{Original references}
\makeatletter
\newlabel{more-morphisms-lemma-separated-locally-quasi-finite-over-affine}{{07S0}{source}{}{sourcepointer-1}{}}
\newlabel{more-morphisms-lemma-separated-locally-quasi-finite-over-affine@cref}{{[originalreference][0][]07S0}{[1][0][]source}{}{}{}}
\makeatother
\crefname{originalremark}{original remark}{original remarks}
\Crefname{originalremark}{Original remark}{Original remarks}
\makeatletter
\newlabel{remark-recall}{{03II}{source}{}{sourcepointer-2}{}}
\newlabel{remark-recall@cref}{{[originalremark][0][]03II}{[1][0][]source}{}{}{}}
\makeatother
\crefname{originalreference}{original reference}{original references}
\Crefname{originalreference}{Original reference}{Original references}
\makeatletter
\newlabel{spaces-morphisms-lemma-composition-quasi-compact}{{03HG}{source}{}{sourcepointer-3}{}}
\newlabel{spaces-morphisms-lemma-composition-quasi-compact@cref}{{[originalreference][0][]03HG}{[1][0][]source}{}{}{}}
\makeatother
\crefname{originalreference}{original reference}{original references}
\Crefname{originalreference}{Original reference}{Original references}
\makeatletter
\newlabel{spaces-morphisms-lemma-diagonal-unramified-morphism}{{05W1}{source}{}{sourcepointer-4}{}}
\newlabel{spaces-morphisms-lemma-diagonal-unramified-morphism@cref}{{[originalreference][0][]05W1}{[1][0][]source}{}{}{}}
\makeatother
\crefname{originalreference}{original reference}{original references}
\Crefname{originalreference}{Original reference}{Original references}
\makeatletter
\newlabel{spaces-morphisms-lemma-etale-universally-injective-open}{{05W5}{source}{}{sourcepointer-5}{}}
\newlabel{spaces-morphisms-lemma-etale-universally-injective-open@cref}{{[originalreference][0][]05W5}{[1][0][]source}{}{}{}}
\makeatother
\crefname{originalreference}{original reference}{original references}
\Crefname{originalreference}{Original reference}{Original references}
\makeatletter
\newlabel{spaces-morphisms-lemma-etale-unramified}{{06CR}{source}{}{sourcepointer-6}{}}
\newlabel{spaces-morphisms-lemma-etale-unramified@cref}{{[originalreference][0][]06CR}{[1][0][]source}{}{}{}}
\makeatother
\crefname{originalreference}{original reference}{original references}
\Crefname{originalreference}{Original reference}{Original references}
\makeatletter
\newlabel{spaces-morphisms-lemma-fibre-product-after-map}{{03KN}{source}{}{sourcepointer-7}{}}
\newlabel{spaces-morphisms-lemma-fibre-product-after-map@cref}{{[originalreference][0][]03KN}{[1][0][]source}{}{}{}}
\makeatother
\crefname{originalreference}{original reference}{original references}
\Crefname{originalreference}{Original reference}{Original references}
\makeatletter
\newlabel{spaces-properties-lemma-etale-local}{{03FS}{source}{}{sourcepointer-8}{}}
\newlabel{spaces-properties-lemma-etale-local@cref}{{[originalreference][0][]03FS}{[1][0][]source}{}{}{}}
\makeatother
\crefname{originalreference}{original reference}{original references}
\Crefname{originalreference}{Original reference}{Original references}
\makeatletter
\newlabel{spaces-properties-lemma-quasi-compact-affine-cover}{{03H6}{source}{}{sourcepointer-9}{}}
\newlabel{spaces-properties-lemma-quasi-compact-affine-cover@cref}{{[originalreference][0][]03H6}{[1][0][]source}{}{}{}}
\makeatother
\crefname{originalreference}{original reference}{original references}
\Crefname{originalreference}{Original reference}{Original references}
\makeatletter
\newlabel{spaces-properties-lemma-quotient-separated}{{0BBM}{source}{}{sourcepointer-10}{}}
\newlabel{spaces-properties-lemma-quotient-separated@cref}{{[originalreference][0][]0BBM}{[1][0][]source}{}{}{}}
\makeatother
\crefname{originalreference}{original reference}{original references}
\Crefname{originalreference}{Original reference}{Original references}
\makeatletter
\newlabel{spaces-properties-proposition-finite-flat-equivalence-global}{{07S6}{source}{}{sourcepointer-11}{}}
\newlabel{spaces-properties-proposition-finite-flat-equivalence-global@cref}{{[originalreference][0][]07S6}{[1][0][]source}{}{}{}}
\makeatother
\title{Stratification of a quasi-compact algebraic space by scheme pieces via symmetric-group quotients}
\author{The Stacks Project Authors}
\date{stacks-project:decent-spaces.tex:local-snapshot-20261012; GNU FDL 1.2 or later}
\begin{document}
\maketitle
\noindent\textit{Editorial note.} No omission markers in-span.\par\medskip
\noindent\textit{Allowed standard prerequisites (editorial).}\begin{itemize}
\item Quasi-compact spaces have affine etale covers. Properties of Spaces, lemma quasi-compact-affine-cover (cross-chapter, named). Provides U -> X surjective etale affine.
\item Diagonals of etale separated morphisms. Morphisms of Spaces, lemmas etale-unramified and diagonal-unramified-morphism (cross-chapter, named). All diagonals are open and closed immersions.
\item Finite sets in affine opens. More on Morphisms, lemma separated-locally-quasi-finite-over-affine (cross-chapter, named). Every finite set of points of W\_p lies in an affine open.
\item Finite flat equivalence relations. Properties of Spaces, proposition finite-flat-equivalence-global (cross-chapter, named). The quotient V\_p = W\_p/S\_p is a scheme.
\item Etale descent of morphism properties and quotients. spaces-properties lemmas etale-local and quotient-separated, and Morphisms of Spaces lemma etale-universally-injective-open (cross-chapter, named). V\_p -> X etale; V\_p separated; Z\_p -> T\_p an open immersion.
\end{itemize}\medskip
\ifdefined\hypertarget\hypertarget{sourcepointer-1}{}\fi
\noindent\textit{Source reference.} Original reference 07S0: focused-seeds/geometry/S9-decent-filter-quasi-compact/ref-tags.txt, source lines 5--5. Navigation pointer only: the counted proof cites this statement by \textbackslash{}ref as a standing input; the deduction the counted proof has to make is carried out inside the counted span, and no part of that deduction is delegated to this pointer.\par\medskip
\ifdefined\hypertarget\hypertarget{sourcepointer-2}{}\fi
\noindent\textit{Source reference.} Original remark 03II: sources/stacks-decent-spaces/pinned-original-source.tex, source lines 459--459. Navigation pointer only: the counted proof cites this statement by \textbackslash{}ref as a standing input; the deduction the counted proof has to make is carried out inside the counted span, and no part of that deduction is delegated to this pointer.\par\medskip
\ifdefined\hypertarget\hypertarget{sourcepointer-3}{}\fi
\noindent\textit{Source reference.} Original reference 03HG: focused-seeds/geometry/S9-decent-filter-quasi-compact/ref-tags.txt, source lines 6--6. Navigation pointer only: the counted proof cites this statement by \textbackslash{}ref as a standing input; the deduction the counted proof has to make is carried out inside the counted span, and no part of that deduction is delegated to this pointer.\par\medskip
\ifdefined\hypertarget\hypertarget{sourcepointer-4}{}\fi
\noindent\textit{Source reference.} Original reference 05W1: focused-seeds/geometry/S9-decent-filter-quasi-compact/ref-tags.txt, source lines 7--7. Navigation pointer only: the counted proof cites this statement by \textbackslash{}ref as a standing input; the deduction the counted proof has to make is carried out inside the counted span, and no part of that deduction is delegated to this pointer.\par\medskip
\ifdefined\hypertarget\hypertarget{sourcepointer-5}{}\fi
\noindent\textit{Source reference.} Original reference 05W5: focused-seeds/geometry/S9-decent-filter-quasi-compact/ref-tags.txt, source lines 8--8. Navigation pointer only: the counted proof cites this statement by \textbackslash{}ref as a standing input; the deduction the counted proof has to make is carried out inside the counted span, and no part of that deduction is delegated to this pointer.\par\medskip
\ifdefined\hypertarget\hypertarget{sourcepointer-6}{}\fi
\noindent\textit{Source reference.} Original reference 06CR: focused-seeds/geometry/S9-decent-filter-quasi-compact/ref-tags.txt, source lines 9--9. Navigation pointer only: the counted proof cites this statement by \textbackslash{}ref as a standing input; the deduction the counted proof has to make is carried out inside the counted span, and no part of that deduction is delegated to this pointer.\par\medskip
\ifdefined\hypertarget\hypertarget{sourcepointer-7}{}\fi
\noindent\textit{Source reference.} Original reference 03KN: focused-seeds/geometry/S9-decent-filter-quasi-compact/ref-tags.txt, source lines 10--10. Navigation pointer only: the counted proof cites this statement by \textbackslash{}ref as a standing input; the deduction the counted proof has to make is carried out inside the counted span, and no part of that deduction is delegated to this pointer.\par\medskip
\ifdefined\hypertarget\hypertarget{sourcepointer-8}{}\fi
\noindent\textit{Source reference.} Original reference 03FS: focused-seeds/geometry/S9-decent-filter-quasi-compact/ref-tags.txt, source lines 11--11. Navigation pointer only: the counted proof cites this statement by \textbackslash{}ref as a standing input; the deduction the counted proof has to make is carried out inside the counted span, and no part of that deduction is delegated to this pointer.\par\medskip
\ifdefined\hypertarget\hypertarget{sourcepointer-9}{}\fi
\noindent\textit{Source reference.} Original reference 03H6: focused-seeds/geometry/S9-decent-filter-quasi-compact/ref-tags.txt, source lines 12--12. Navigation pointer only: the counted proof cites this statement by \textbackslash{}ref as a standing input; the deduction the counted proof has to make is carried out inside the counted span, and no part of that deduction is delegated to this pointer.\par\medskip
\ifdefined\hypertarget\hypertarget{sourcepointer-10}{}\fi
\noindent\textit{Source reference.} Original reference 0BBM: focused-seeds/geometry/S9-decent-filter-quasi-compact/ref-tags.txt, source lines 13--13. Navigation pointer only: the counted proof cites this statement by \textbackslash{}ref as a standing input; the deduction the counted proof has to make is carried out inside the counted span, and no part of that deduction is delegated to this pointer.\par\medskip
\ifdefined\hypertarget\hypertarget{sourcepointer-11}{}\fi
\noindent\textit{Source reference.} Original reference 07S6: focused-seeds/geometry/S9-decent-filter-quasi-compact/ref-tags.txt, source lines 14--14. Navigation pointer only: the counted proof cites this statement by \textbackslash{}ref as a standing input; the deduction the counted proof has to make is carried out inside the counted span, and no part of that deduction is delegated to this pointer.\par\medskip
\begin{lemma}
\label{lemma-filter-quasi-compact}
Let $S$ be a scheme. Let $X$ be a quasi-compact algebraic space
over $S$. There exist open subspaces
$$
\ldots \subset U_4 \subset U_3 \subset U_2 \subset U_1 = X
$$
with the following properties:
\begin{enumerate}
\item setting $T_p = U_p \setminus U_{p + 1}$ (with reduced induced subspace
structure) there exists a separated scheme $V_p$ and a surjective \'etale
morphism $f_p : V_p \to U_p$ such that $f_p^{-1}(T_p) \to T_p$ is an
isomorphism,
\item if $x \in |X|$ can be represented by a quasi-compact morphism
$\Spec(k) \to X$ from a field, then $x \in T_p$ for some $p$.
\end{enumerate}
\end{lemma}

\begin{proof}
By Properties of Spaces, Lemma
\ref{spaces-properties-lemma-quasi-compact-affine-cover}
we can choose an affine scheme $U$ and a surjective \'etale morphism
$U \to X$. For $p \geq 0$ set
$$
W_p = U \times_X \ldots \times_X U \setminus \text{all diagonals}
$$
where the fibre product has $p$ factors. Since $U$ is separated,
the morphism $U \to X$ is separated and all fibre products
$U \times_X \ldots \times_X U$ are separated schemes. Since $U \to X$ is
separated the diagonal $U \to U \times_X U$ is a closed immersion. Since
$U \to X$ is \'etale the diagonal $U \to U \times_X U$ is an open
immersion, see Morphisms of Spaces, Lemmas
\ref{spaces-morphisms-lemma-etale-unramified} and
\ref{spaces-morphisms-lemma-diagonal-unramified-morphism}.
Similarly, all the diagonal morphisms are open and closed immersions and
$W_p$ is an open and closed subscheme of $U \times_X \ldots \times_X U$.
Moreover, the morphism
$$
U \times_X \ldots \times_X U \longrightarrow
U \times_{\Spec(\mathbf{Z})} \ldots \times_{\Spec(\mathbf{Z})} U
$$
is locally quasi-finite and separated (Morphisms of Spaces,
Lemma \ref{spaces-morphisms-lemma-fibre-product-after-map})
and its target is an affine scheme. Hence every finite set of points of
$U \times_X \ldots \times_X U$ is contained in an affine open, see
More on Morphisms, Lemma
\ref{more-morphisms-lemma-separated-locally-quasi-finite-over-affine}.
Therefore, the same is true for $W_p$.
There is a free action of the symmetric group $S_p$ on $W_p$ over $X$
(because we threw out the fix point locus from
$U \times_X \ldots \times_X U$). By the above and
Properties of Spaces, Proposition
\ref{spaces-properties-proposition-finite-flat-equivalence-global}
the quotient $V_p = W_p/S_p$ is a scheme. Since the action of
$S_p$ on $W_p$ was over $X$, there is a morphism $V_p \to X$.
Since $W_p \to X$ is \'etale and since $W_p \to V_p$ is surjective
\'etale, it follows that also $V_p \to X$ is \'etale, see
Properties of Spaces, Lemma \ref{spaces-properties-lemma-etale-local}.
Observe that $V_p$ is a separated scheme by
Properties of Spaces, Lemma
\ref{spaces-properties-lemma-quotient-separated}.

\medskip\noindent
We let $U_p \subset X$ be the open subspace which is the
image of $V_p \to X$. By construction a morphism $\Spec(k) \to X$ with
$k$ algebraically closed, factors through $U_p$ if and only if
$U \times_X \Spec(k)$ has $\geq p$ points; as usual observe that
$U \times_X \Spec(k)$ is scheme theoretically a disjoint union of
(possibly infinitely many) copies of $\Spec(k)$, see
Remark \ref{remark-recall}. It follows that
the $U_p$ give a filtration of $X$ as stated in the lemma.
Moreover, our morphism $\Spec(k) \to X$ factors through $T_p$
if and only if $U \times_X \Spec(k)$ has exactly $p$ points.
In this case we see that $V_p \times_X \Spec(k)$ has exactly one point.
Set $Z_p = f_p^{-1}(T_p) \subset V_p$. This is a closed subscheme of $V_p$.
Then $Z_p \to T_p$ is an \'etale morphism between
algebraic spaces which induces a bijection on $k$-valued
points for any algebraically closed field $k$. To be sure this
implies that $Z_p \to T_p$ is universally injective, whence an
open immersion by
Morphisms of Spaces, Lemma
\ref{spaces-morphisms-lemma-etale-universally-injective-open}
hence an isomorphism and (1) has been proved.

\medskip\noindent
Let $x : \Spec(k) \to X$ be a quasi-compact morphism where $k$ is a field.
Then the composition $\Spec(\overline{k}) \to \Spec(k) \to X$ is quasi-compact
as well (Morphisms of Spaces, Lemma
\ref{spaces-morphisms-lemma-composition-quasi-compact}).
In this case the scheme $U \times_X \Spec(\overline{k})$ is
quasi-compact. In view of the fact (seen above) that it is a disjoint union
of copies of $\Spec(\overline{k})$ we find that it has finitely many points.
If the number of points is $p$, then we see that indeed $x \in T_p$ and
the proof is finished.
\end{proof}

\end{document}
